\documentclass[a4paper, 12pt]{article}

\usepackage[portuges]{babel}
\usepackage[utf8]{inputenc}
\usepackage{amsmath}
\usepackage{amsthm}
\usepackage{indentfirst}
\usepackage{graphicx}
\usepackage[colorinlistoftodos]{todonotes}
\usepackage{verbatim}
\usepackage{textcomp}
\usepackage{gensymb}
\usepackage{float}
\usepackage{url}
\usepackage{listings}
% Cor do texto nas tabelas: verde = passo sem problema, vermelho = erro.
\newcommand{\ok}[1]{\textcolor{green!45!black}{#1}}
\newcommand{\erro}[1]{\textcolor{red!70!black}{#1}}
% Margens no padrão ABNT (3 cm em cima/esquerda, 2 cm embaixo/direita).
% As margens padrão do article são bem maiores e não caberia tudo em 4 páginas.
\usepackage[top=3cm, left=3cm, bottom=2cm, right=2cm]{geometry}
% Listas (itemize/enumerate) mais compactas, pra economizar espaço.
\usepackage{enumitem}
\setlist{itemsep=1pt, topsep=3pt}

% Ambiente "Lema" (o \begin{proof} ... \end{proof} já vem do amsthm)
\newtheorem{lema}{Lema}

% Configuração dos códigos em C.
% O "literate" serve pra os acentos dos comentários em português não quebrarem o listings.
\lstset{
    language=C,
    basicstyle=\ttfamily\footnotesize,
    keywordstyle=\bfseries\color{blue!70!black},
    commentstyle=\itshape\color{green!40!black},
    stringstyle=\color{red!60!black},
    numbers=left,
    numberstyle=\tiny\color{gray},
    frame=single,
    breaklines=true,
    showstringspaces=false,
    tabsize=4,
    literate={á}{{\'a}}1 {à}{{\`a}}1 {ã}{{\~a}}1 {â}{{\^a}}1
             {é}{{\'e}}1 {ê}{{\^e}}1 {í}{{\'\i}}1
             {ó}{{\'o}}1 {õ}{{\~o}}1 {ô}{{\^o}}1 {ú}{{\'u}}1 {ç}{{\c c}}1
             {Á}{{\'A}}1 {É}{{\'E}}1 {Í}{{\'I}}1 {Ó}{{\'O}}1 {Ú}{{\'U}}1
             {Ã}{{\~A}}1 {Õ}{{\~O}}1 {Ç}{{\c C}}1
}

\begin{document}

\begin{titlepage}
\begin{center}
\textbf{\LARGE Universidade de Brasília}\\[0.5cm]
\textbf{\large Departamento de Ciência da Computação}\\[0.2cm]
\vspace{20pt}

\par
\vspace{32pt}
% O \par DENTRO das chaves faz o LaTeX usar o espaçamento de linha do \LARGE
% (antes ele usava o do tamanho normal, por isso as linhas ficavam grudadas).
{\LARGE\bfseries Relatório do Trabalho 1 de\\[0.25cm] Programação Concorrente\par}
\vspace{30pt}
\Large {Evelyn Soares Pereira	17/0102785}\\[0.1cm]

\end{center}

\par
\vfill
\begin{center}
{{\normalsize Brasília}\\
{\normalsize \today}}
\end{center}
\end{titlepage}

\newpage

% Limite: 4 páginas, sem contar capa e códigos inseridos no texto.

%======================================================================
\section{Introdução}
%======================================================================

Em programação concorrente, \emph{threads} que cooperam por memória
compartilhada correm o risco de \emph{condições de corrida}: se duas delas
acessam o mesmo dado sem coordenação e pelo menos uma escreve, o resultado
passa a depender da ordem em que o escalonador intercala suas instruções.
Como essa ordem muda a cada execução, os erros aparecem em uma execução e
podem sumir na seguinte. Para evitá-los, as \emph{threads} são
sincronizadas com mecanismos como \emph{locks}, variáveis de condição e
semáforos.

Este trabalho aplica esses mecanismos a um problema de Manna e Pnueli,
apresentado por Ben-Ari (2006, Apêndice C, exercício 8): calcular o
coeficiente binomial $\binom{n}{k}$ com dois processos, um para o numerador e
outro para o denominador. O denominador entrega resultados parciais ao
numerador, que divide assim que os recebe, impedindo que os valores
intermediários cresçam demais. Essa comunicação acontece por memória
compartilhada e, sem sincronização, está sujeita a condições de corrida. A
solução foi implementada em C com POSIX Pthreads. A
Seção~\ref{sec:formalizacao} formaliza o problema, a
Seção~\ref{sec:algoritmo} descreve o algoritmo, a implementação e os testes,
e a Seção~\ref{sec:conclusao} conclui o trabalho.


%======================================================================
\section{Formalização do Problema Proposto}
\label{sec:formalizacao}
%======================================================================

\subsection{O Coeficiente Binomial}

Para inteiros $0 \le k \le n$,
\begin{equation}
    \binom{n}{k} = \frac{n!}{(n-k)!\,k!}
                 = \frac{n \cdot (n-1) \cdots (n-k+1)}{1 \cdot 2 \cdots k},
\end{equation}
ou seja, há $k$ fatores no numerador e $k$ no denominador. Calcular o
numerador inteiro e só dividir no final é inviável mesmo para valores
modestos: em $\binom{30}{15}$, o numerador
$30 \cdot 29 \cdots 16 \approx 2{,}0 \times 10^{20}$ não cabe em 64 bits
($\approx 1{,}8 \times 10^{19}$), embora o resultado ($155\,117\,520$)
caiba em 32. A proposta do exercício é intercalar as operações: no passo
$i$, multiplicar o valor acumulado pelo $i$-ésimo fator do numerador e
dividir pelo $i$-ésimo fator do denominador (Tabela~\ref{tab:exemplo}).

\begin{table}[H]
\centering
\caption{Cálculo de $\binom{10}{3}$ intercalando multiplicações e divisões.}
\label{tab:exemplo}
\begin{tabular}{ccccc}
\hline
Passo $i$ & Multiplica por & Antes da divisão & Divide por & Depois da divisão \\
\hline
1 & 10 & 10  & 1 & 10  \\
2 & 9  & 90  & 2 & 45  \\
3 & 8  & 360 & 3 & 120 \\
\hline
\end{tabular}
\end{table}

Como a divisão inteira em C descarta o resto sem aviso, esse método só
funciona se toda divisão intermediária for exata, o que é provado a seguir.

\subsection{Divisibilidade dos Resultados Parciais}

Chamamos de $P(i,j)$ o produto de $i$ inteiros consecutivos começando em
$j$:
\[
    P(i,j) = j \cdot (j+1) \cdots (j+i-1).
    \qquad \text{Por exemplo, } P(3,5) = 5 \cdot 6 \cdot 7 = 210.
\]

\begin{lema}\label{lema:div}
    Para todos os inteiros $i \ge 1$ e $j \ge 1$, $i!$ divide $P(i,j)$.
\end{lema}

A demonstração usa uma identidade auxiliar. Para $i \ge 2$, os produtos
$P(i,j)$ e $P(i,j+1)$ têm em comum os fatores $(j+1) \cdots (j+i-1)$, que
formam $P(i-1,\,j+1)$. Assim,
\[
    P(i,j) = j \cdot P(i-1,\,j+1)
    \qquad \text{e} \qquad
    P(i,j+1) = P(i-1,\,j+1) \cdot (j+i),
\]
e, subtraindo uma da outra,
\begin{equation}\label{eq:identidade}
    P(i,j+1) - P(i,j) = \big[(j+i) - j\big] \cdot P(i-1,\,j+1)
                      = i \cdot P(i-1,\,j+1).
\end{equation}
Por exemplo, $P(3,6) - P(3,5) = 6 \cdot 7 \cdot 8 - 5 \cdot 6 \cdot 7
= (8-5) \cdot 6 \cdot 7 = 3 \cdot P(2,6)$.

\begin{proof}
Por indução em $i$.

\emph{Base ($i = 1$):} $P(1,j) = j$, e $1! = 1$ divide qualquer inteiro.

\emph{Passo:} seja $i \ge 2$ e suponha, como hipótese de indução, que
$(i-1)!$ divide $P(i-1,m)$ para todo $m \ge 1$. Vamos mostrar, agora por
indução em $j$, que $i!$ divide $P(i,j)$ para todo $j \ge 1$.
\begin{itemize}
    \item $j = 1$: $P(i,1) = 1 \cdot 2 \cdots i = i!$, que é divisível
          por $i!$.
    \item De $j$ para $j+1$: suponha que $i!$ divide $P(i,j)$. Pela
          identidade~\eqref{eq:identidade},
          \[ P(i,j+1) = P(i,j) + i \cdot P(i-1,\,j+1). \]
          A primeira parcela é múltipla de $i!$ pela suposição. Na segunda,
          a hipótese de indução (com $m = j+1$) dá
          $P(i-1,\,j+1) = c \cdot (i-1)!$ para algum inteiro $c$, logo
          $i \cdot P(i-1,\,j+1) = c \cdot i \cdot (i-1)! = c \cdot i!$.
          A soma de dois múltiplos de $i!$ é múltipla de $i!$. \qedhere
\end{itemize}
\end{proof}

\paragraph{Aplicação ao problema.} Depois de $i$ multiplicações, o
numerador acumulou $N_i = n \cdot (n-1) \cdots (n-i+1)$. Lido de trás para a
frente, esse é um produto de $i$ inteiros consecutivos começando em
$n-i+1$, ou seja, $N_i = P(i,\,n-i+1)$; por exemplo,
$10 \cdot 9 \cdot 8 = 8 \cdot 9 \cdot 10 = P(3,8)$. Como $i \le k \le n$, o
início $n-i+1$ é pelo menos 1, e o lema garante que $i!$ divide $N_i$.

Isso explica a Tabela~\ref{tab:exemplo}. Antes da $i$-ésima divisão, o
acumulado já foi dividido por $1, 2, \dots, i-1$, ou seja, por $(i-1)!$.
Então
\[
    \underbrace{\frac{N_i}{(i-1)!}}_{\text{antes da divisão}} \div\; i
    \;=\; \frac{N_i}{i!} \;=\; \binom{n}{i},
\]
que é inteiro pelo lema. No passo 3 da tabela, por exemplo,
$720 / 2! = 360$ e $360 / 3 = 120 = 720 / 3!$. Por isso o denominador não
precisa enviar $i!$: basta enviar o próximo fator $i$, e as divisões
sucessivas por $1, 2, \dots, i$ têm o mesmo efeito que dividir por $i!$. Em
$\binom{30}{15}$, o maior valor acumulado fica em torno de
$2{,}3 \times 10^{9}$, contra $2{,}0 \times 10^{20}$ do numerador completo.


\subsection{Processos, Dados Compartilhados e Condições de Corrida}

O programa tem três \emph{threads}:
\begin{itemize}
    \item \textbf{Denominador} ($T_d$): gera, em ordem, os fatores
          $1, 2, \dots, k$ e os escreve num \emph{buffer} compartilhado;
    \item \textbf{Numerador} ($T_n$): para $i = 1, \dots, k$, multiplica o
          seu acumulado por $n-i+1$, lê o próximo fator do \emph{buffer} e
          divide o acumulado por ele;
    \item \textbf{Principal}: lê $n$ e $k$, cria $T_d$ e $T_n$, espera as
          duas terminarem (\texttt{pthread\_join}) e imprime o resultado.
\end{itemize}

O acumulado pertence só a $T_n$ (a principal só o lê depois do
\texttt{pthread\_join}), então não há disputa por ele. O que é compartilhado
é o \emph{buffer}: as posições de memória onde $T_d$ escreve os fatores e de
onde $T_n$ os lê, junto com as variáveis que indicam quantos fatores estão
disponíveis. Sem sincronização, ele está sujeito a três condições de
corrida:
\begin{enumerate}
    \item \textbf{Leitura antes da escrita:} $T_n$ lê uma posição em que
          $T_d$ ainda não escreveu e obtém um valor antigo ou o valor
          inicial (se for 0, divide por zero).
    \item \textbf{Sobrescrita:} $T_d$ escreve um fator por cima de outro que
          $T_n$ ainda não leu, e esse divisor se perde.
    \item \textbf{Atualização perdida:} se as duas \emph{threads} alteram a
          mesma variável de controle (por exemplo, \texttt{qtd++} em $T_d$ e
          \texttt{qtd-{}-} em $T_n$), cada operação é feita em três passos
          (ler, alterar, escrever), e uma alteração pode apagar a outra.
\end{enumerate}

A Tabela~\ref{tab:corrida} mostra uma execução de $\binom{5}{2} = 10$ sem
sincronização, com um \emph{buffer} de uma posição, em que as corridas 1 e 2
acontecem e o programa imprime 4 sem nenhum erro aparente.

\begin{table}[H]
\centering
\caption{Execução sem sincronização de $\binom{5}{2}$, passo a passo.}
\label{tab:corrida}
\begin{tabular}{clclcl}
\hline
Passo & $T_d$ & \emph{buffer} & $T_n$ & Acumulado & Problema \\
\hline
\ok{1}   & \ok{--}          & \ok{--}  & \ok{multiplica por 5}        & \ok{5}   & \\
\ok{2}   & \ok{escreve 1}   & \ok{1}   & \ok{--}                      & \ok{5}   & \\
\erro{3} & \erro{escreve 2} & \erro{2} & \erro{--}                    & \erro{5} & \erro{o 1 se perdeu (corrida 2)} \\
\erro{4} & \erro{--}        & \erro{2} & \erro{lê 2 e divide: $5/2$}  & \erro{2} & \erro{divisão truncada} \\
\ok{5}   & \ok{--}          & \ok{2}   & \ok{multiplica por 4}        & \ok{8}   & \\
\erro{6} & \erro{--}        & \erro{2} & \erro{lê 2 de novo: $8/2$}   & \erro{\textbf{4}} & \erro{valor antigo (corrida 1)} \\
\hline
\end{tabular}
\end{table}

Uma solução correta deve garantir:
\begin{itemize}
    \item \textbf{entrega exata:} cada fator é lido por $T_n$ exatamente uma
          vez, na ordem em que foi escrito;
    \item \textbf{exclusão mútua} no acesso às variáveis do \emph{buffer};
    \item \textbf{término sem \emph{deadlock}:} as duas \emph{threads}
          terminam após $k$ iterações.
\end{itemize}
Com as duas primeiras, $T_n$ repete exatamente a Tabela~\ref{tab:exemplo}
e, pelo Lema~\ref{lema:div}, obtém $\binom{n}{k}$. Também é desejável que as
esperas sejam por bloqueio, e não por espera ocupada.


%======================================================================
\section{Descrição do Algoritmo Desenvolvido para Solução do Problema Proposto}
\label{sec:algoritmo}
%======================================================================

\subsection{Visão Geral}
% A ideia da solução em texto. Uma figura da interação entre as threads ajuda.


\subsection{Mecanismos de Sincronização Utilizados}
% Quais mecanismos (locks, variáveis de condição, semáforos...) e POR QUE.


\subsection{Pseudocódigo}
% Pseudocódigo de cada thread.
% \begin{lstlisting}[language={}, numbers=none]
% ...
% \end{lstlisting}


\subsection{Corretude}
% Argumentar por que a solução funciona: resultado correto,
% exclusão mútua, ausência de deadlock e de starvation.


\subsection{Implementação}
% - Arquivo(s), como compilar e executar:
%     gcc -o NOME NOME.c -lpthread
% - Trechos importantes do código:
% \lstinputlisting[caption={Descrição do trecho}]{NOME.c}


\subsection{Testes e Resultados}
% Quais valores de n e k foram testados, saída esperada x obtida.
% \begin{verbatim}
% saída do programa aqui
% \end{verbatim}


%======================================================================
\section{Conclusão}
\label{sec:conclusao}
%======================================================================
% - O que foi feito.
% - Dificuldades e o que aprendeu.
% - Limitações e possíveis extensões.


%======================================================================
\section{Referências}
%======================================================================

\sloppy
\noindent
BEN-ARI, Mordechai. \textbf{Principles of Concurrent and Distributed Programming}. 2. ed. Harlow: Addison-Wesley, 2006. ISBN 978-0-321-31283-9.

% \vspace{0.4cm}
% \noindent
% ALCHIERI, Eduardo A. P. \textbf{Programação Concorrente}. Notas de aula. Brasília: Universidade de Brasília, 2026.

% \vspace{0.4cm}
% \noindent
% (Outras: documentação do POSIX Threads, etc.)

\fussy

\end{document}
